% Options for packages loaded elsewhere
\PassOptionsToPackage{unicode}{hyperref}
\PassOptionsToPackage{hyphens}{url}
\PassOptionsToPackage{dvipsnames,svgnames,x11names}{xcolor}
%
\documentclass[
  12pt,
]{article}
\usepackage{amsmath,amssymb}
\usepackage{iftex}
\ifPDFTeX
  \usepackage[T1]{fontenc}
  \usepackage[utf8]{inputenc}
  \usepackage{textcomp} % provide euro and other symbols
\else % if luatex or xetex
  \usepackage{unicode-math} % this also loads fontspec
  \defaultfontfeatures{Scale=MatchLowercase}
  \defaultfontfeatures[\rmfamily]{Ligatures=TeX,Scale=1}
\fi
\usepackage{lmodern}
\ifPDFTeX\else
  % xetex/luatex font selection
\fi
% Use upquote if available, for straight quotes in verbatim environments
\IfFileExists{upquote.sty}{\usepackage{upquote}}{}
\IfFileExists{microtype.sty}{% use microtype if available
  \usepackage[]{microtype}
  \UseMicrotypeSet[protrusion]{basicmath} % disable protrusion for tt fonts
}{}
\makeatletter
\@ifundefined{KOMAClassName}{% if non-KOMA class
  \IfFileExists{parskip.sty}{%
    \usepackage{parskip}
  }{% else
    \setlength{\parindent}{0pt}
    \setlength{\parskip}{6pt plus 2pt minus 1pt}}
}{% if KOMA class
  \KOMAoptions{parskip=half}}
\makeatother
\usepackage{xcolor}
\usepackage[margin=1in]{geometry}
\usepackage{color}
\usepackage{fancyvrb}
\newcommand{\VerbBar}{|}
\newcommand{\VERB}{\Verb[commandchars=\\\{\}]}
\DefineVerbatimEnvironment{Highlighting}{Verbatim}{commandchars=\\\{\}}
% Add ',fontsize=\small' for more characters per line
\newenvironment{Shaded}{}{}
\newcommand{\AlertTok}[1]{\textcolor[rgb]{1.00,0.00,0.00}{\textbf{#1}}}
\newcommand{\AnnotationTok}[1]{\textcolor[rgb]{0.38,0.63,0.69}{\textbf{\textit{#1}}}}
\newcommand{\AttributeTok}[1]{\textcolor[rgb]{0.49,0.56,0.16}{#1}}
\newcommand{\BaseNTok}[1]{\textcolor[rgb]{0.25,0.63,0.44}{#1}}
\newcommand{\BuiltInTok}[1]{\textcolor[rgb]{0.00,0.50,0.00}{#1}}
\newcommand{\CharTok}[1]{\textcolor[rgb]{0.25,0.44,0.63}{#1}}
\newcommand{\CommentTok}[1]{\textcolor[rgb]{0.38,0.63,0.69}{\textit{#1}}}
\newcommand{\CommentVarTok}[1]{\textcolor[rgb]{0.38,0.63,0.69}{\textbf{\textit{#1}}}}
\newcommand{\ConstantTok}[1]{\textcolor[rgb]{0.53,0.00,0.00}{#1}}
\newcommand{\ControlFlowTok}[1]{\textcolor[rgb]{0.00,0.44,0.13}{\textbf{#1}}}
\newcommand{\DataTypeTok}[1]{\textcolor[rgb]{0.56,0.13,0.00}{#1}}
\newcommand{\DecValTok}[1]{\textcolor[rgb]{0.25,0.63,0.44}{#1}}
\newcommand{\DocumentationTok}[1]{\textcolor[rgb]{0.73,0.13,0.13}{\textit{#1}}}
\newcommand{\ErrorTok}[1]{\textcolor[rgb]{1.00,0.00,0.00}{\textbf{#1}}}
\newcommand{\ExtensionTok}[1]{#1}
\newcommand{\FloatTok}[1]{\textcolor[rgb]{0.25,0.63,0.44}{#1}}
\newcommand{\FunctionTok}[1]{\textcolor[rgb]{0.02,0.16,0.49}{#1}}
\newcommand{\ImportTok}[1]{\textcolor[rgb]{0.00,0.50,0.00}{\textbf{#1}}}
\newcommand{\InformationTok}[1]{\textcolor[rgb]{0.38,0.63,0.69}{\textbf{\textit{#1}}}}
\newcommand{\KeywordTok}[1]{\textcolor[rgb]{0.00,0.44,0.13}{\textbf{#1}}}
\newcommand{\NormalTok}[1]{#1}
\newcommand{\OperatorTok}[1]{\textcolor[rgb]{0.40,0.40,0.40}{#1}}
\newcommand{\OtherTok}[1]{\textcolor[rgb]{0.00,0.44,0.13}{#1}}
\newcommand{\PreprocessorTok}[1]{\textcolor[rgb]{0.74,0.48,0.00}{#1}}
\newcommand{\RegionMarkerTok}[1]{#1}
\newcommand{\SpecialCharTok}[1]{\textcolor[rgb]{0.25,0.44,0.63}{#1}}
\newcommand{\SpecialStringTok}[1]{\textcolor[rgb]{0.73,0.40,0.53}{#1}}
\newcommand{\StringTok}[1]{\textcolor[rgb]{0.25,0.44,0.63}{#1}}
\newcommand{\VariableTok}[1]{\textcolor[rgb]{0.10,0.09,0.49}{#1}}
\newcommand{\VerbatimStringTok}[1]{\textcolor[rgb]{0.25,0.44,0.63}{#1}}
\newcommand{\WarningTok}[1]{\textcolor[rgb]{0.38,0.63,0.69}{\textbf{\textit{#1}}}}
\usepackage{longtable,booktabs,array}
\usepackage{calc} % for calculating minipage widths
% Correct order of tables after \paragraph or \subparagraph
\usepackage{etoolbox}
\makeatletter
\patchcmd\longtable{\par}{\if@noskipsec\mbox{}\fi\par}{}{}
\makeatother
% Allow footnotes in longtable head/foot
\IfFileExists{footnotehyper.sty}{\usepackage{footnotehyper}}{\usepackage{footnote}}
\makesavenoteenv{longtable}
\setlength{\emergencystretch}{3em} % prevent overfull lines
\providecommand{\tightlist}{%
  \setlength{\itemsep}{0pt}\setlength{\parskip}{0pt}}
\setcounter{secnumdepth}{-\maxdimen} % remove section numbering
\ifLuaTeX
  \usepackage{selnolig}  % disable illegal ligatures
\fi
\IfFileExists{bookmark.sty}{\usepackage{bookmark}}{\usepackage{hyperref}}
\IfFileExists{xurl.sty}{\usepackage{xurl}}{} % add URL line breaks if available
\urlstyle{same}
\hypersetup{
  pdftitle={GPU Sharing Mechanisms for Multi-Tenant Inference Serving},
  pdfauthor={Operating Systems CEP --- shared candidate for two independent groups},
  colorlinks=true,
  linkcolor={Maroon},
  filecolor={Maroon},
  citecolor={Blue},
  urlcolor={Blue},
  pdfcreator={LaTeX via pandoc}}

\title{GPU Sharing Mechanisms for Multi-Tenant Inference Serving}
\usepackage{etoolbox}
\makeatletter
\providecommand{\subtitle}[1]{% add subtitle to \maketitle
  \apptocmd{\@title}{\par {\large #1 \par}}{}{}
}
\makeatother
\subtitle{Systematic Literature Review Protocol v1.2 (Proposed
Pre-Search Draft)}
\author{Operating Systems CEP --- shared candidate for two independent
groups}
\date{27 September 2026}

\begin{document}
\maketitle

\textbf{Version:} v1.2, proposed pre-search revision, 27 September 2026
(PKT)\\
\textbf{Status:} Not frozen. The four database interfaces, exact
executable queries, access, seed recall, group ownership and reviewer
sign-off remain unverified. No formal search, screened study or PRISMA
count is claimed.\\
\textbf{Ownership:} Shared candidate for two independent three-student
groups. Each group must adapt, pilot, date and freeze its own protocol
before formal selection. Neither may inherit the other\textquotesingle s
decisions or counts. The root v1.0/v1.1 history remains in
\texttt{evidence/protocol\_amendments.md}.

The twelve sections below follow the instructor\textquotesingle s
requested order. PRISMA-P informs protocol transparency {[}1{]}, PRISMA
2020 informs final review reporting {[}2{]}, PRISMA-S informs
reproducible search documentation {[}3{]}, and software-engineering SLR
guidance informs the review process {[}4{]}. These reporting guidelines
are \textbf{method references}, not included primary studies.

\hypertarget{1-research-objective}{%
\section{1. Research objective}\label{1-research-objective}}

Identify, critically appraise and synthesize peer-reviewed primary
studies evaluating how \textbf{one physical GPU on a single node} is
shared by multiple inference tenants. Compare temporal time-slicing,
CUDA Multi-Process Service (MPS), Multi-Instance GPU (MIG), CUDA/GPU API
remoting or virtualization, and hybrids. The intended output is a
conditional engineering choice: which mechanism fits a specified model
and arrival pattern, GPU generation/configuration, latency service-level
objective (SLO), utilization or tenant-density goal, and isolation
requirement. Treat performance interference, resource/memory separation,
security isolation and fault containment as distinct observations; no
mechanism name proves a measured outcome {[}10{]}, {[}11{]}.

\hypertarget{2-research-questions}{%
\section{2. Research questions}\label{2-research-questions}}

\begin{enumerate}
\def\labelenumi{\arabic{enumi}.}
\tightlist
\item
  \textbf{RQ1 (alternatives):} Which of the four assigned mechanism
  families and hybrids have been evaluated for concurrent inference
  tenants on a single GPU, and what control/resource boundaries were
  evaluated?
\item
  \textbf{RQ2 (measurement):} How are utilization, throughput/goodput,
  tenant density, p50/p95/p99 latency, SLO violations, interference and
  isolation measured? Which studies can actually be compared?
\item
  \textbf{RQ3 (conditions):} How do workload/model, request arrival and
  batching, tenant mix, GPU generation, MPS settings, MIG profile and
  remoting path affect the reported outcomes?
\item
  \textbf{RQ4 (trade-offs):} What performance, isolation, implementation
  and operational trade-offs, conflicting findings and evidence gaps are
  reported?
\item
  \textbf{RQ5 (decision):} Under stated SLO, utilization/density,
  isolation and hardware conditions, what mechanism or combination is
  supported, and how strong is that support?
\end{enumerate}

The group must map each final RQ to the extraction fields in §11 and
synthesis outputs in §12 before freezing.

\hypertarget{3-selected-academic-databases}{%
\section{3. Selected academic
databases}\label{3-selected-academic-databases}}

The \textbf{proposed formal sources} are IEEE Xplore
(engineering/systems), ACM Digital Library (computing/systems),
ScienceDirect (Elsevier engineering/computing journals), and Springer
Nature Link (computing journals/proceedings). NED\textquotesingle s
E-Resources page lists these platforms in its E-Databases area {[}5{]}.
Verify each group\textquotesingle s actual entitlement, indexed
collection and export route before freeze; a publisher collection is not
an exhaustive cross-publisher index. If unavailable, document a
justified replacement and revised coverage.

Scopus and Web of Science, if accessible, supplement cross-publisher
coverage/citation checking; any systematic run must be logged and
screened like a formal source. Google Scholar and backward/forward
references are supplementary discovery, logged separately. Check whether
USENIX and other systems proceedings are missed by the proposed
publisher platforms. Do not mistake publisher access for database
coverage.

\hypertarget{4-search-keywords}{%
\section{4. Search keywords}\label{4-search-keywords}}

\begin{longtable}[]{@{}p{0.23\textwidth}p{0.70\textwidth}@{}}
\toprule\noalign{}
Concept & Candidate keywords and spelling variants \\
\midrule\noalign{}
\endhead
\bottomrule\noalign{}
\endlastfoot
Hardware & GPU; graphics processing unit; accelerator \\
Inference & inference; model serving; deep learning; DNN; neural
network; LLM; online serving \\
Temporal sharing & time slicing; time-slicing; time sharing; context
switching; preemptive scheduling \\
MPS & MPS; multi-process service; CUDA MPS; process concurrency \\
MIG/partitioning & MIG; multi-instance GPU; GPU partitioning; spatial
sharing \\
API remoting & API remoting; CUDA remoting; GPU proxy; API interception;
virtual GPU; GPU virtualization \\
Outcomes/setting (optional refinements) & multi-tenant; co-location;
utilization; latency; tail latency; QoS; SLO; interference; isolation \\
\end{longtable}

The mandatory query logic is hardware AND inference AND an OR of
\textbf{all four} mechanism families. Do not require
\texttt{multi-tenant} or an outcome word in every record: relevant seed
papers may lack that literal term {[}12{]}. Pilot spelling,
tokenization, wildcard and subject-field choices for each interface
before freeze. MPS and MIG alone can have unrelated meanings, so pair
with GPU/inference and inspect precision.

\hypertarget{5-search-strings}{%
\section{5. Search strings}\label{5-search-strings}}

\textbf{Candidate Boolean logic (not yet an executed search):}

\begin{Shaded}
\begin{Highlighting}[]
\NormalTok{(GPU OR "graphics processing unit")}
\NormalTok{AND (inference OR "model serving" OR "deep learning" OR DNN OR LLM)}
\NormalTok{AND ("time slicing" OR "time sharing" OR "context switching"}
\NormalTok{     OR MPS OR "multi{-}process service" OR MIG OR "multi{-}instance GPU"}
\NormalTok{     OR "GPU partitioning" OR "API remoting" OR "CUDA remoting"}
\NormalTok{     OR "GPU proxy" OR "GPU virtualization")}
\end{Highlighting}
\end{Shaded}

\begin{longtable}[]{@{}p{0.23\textwidth}p{0.70\textwidth}@{}}
\toprule\noalign{}
Interface to pilot & Proposed entry and mandatory validation \\
\midrule\noalign{}
\endhead
\bottomrule\noalign{}
\endlastfoot
IEEE Xplore Command Search & Enter the candidate logic; verify uppercase
Boolean, phrase grouping, metadata field, date/type facets and displayed
executed query. \\
ACM DL Advanced Search & Enter concepts in the builder in the
\textbf{ACM full-text collection}; capture its generated syntax, field
scope and collection (not automatically the broader ACM Guide). \\
ScienceDirect Advanced Search & Test the broad logic in the Advanced
Search terms field; record accepted syntax, whether full text or
title/abstract/keywords was searched, and applied article/date
filters. \\
Springer Nature Link Advanced Search & Test grouping, phrases, uppercase
Boolean, field scope and year/content-type facets; record the executed
representation. \\
\end{longtable}

If clause limits or seed misses require it, run \textbf{four separately
logged queries per platform} with the common GPU AND inference blocks
and one mechanism-family OR block from §4 each; deduplicate their
exported records. A family split or field change affecting recall must
be chosen and recorded \textbf{before} formal freeze.
\texttt{evidence/query\_validation.md} contains the existing unexecuted
platform inputs; each group must replace proposals with its actual
accepted query, field, URL, filters, pilot hit count and seed-retrieval
result. Suggested seed \emph{candidates} span PipeSwitch {[}12{]},
ParvaGPU {[}13{]} and Torpor {[}14{]}; none is pre-included. A seed
absent from its publisher collection tests coverage, not necessarily
query syntax.

For every later \textbf{formal} run log group ID, database/collection,
exact executed query, fields, filters, date/time, URL, total hits,
export filename and reviewer. Keep pilot and formal runs separate.
PRISMA-S {[}3{]} motivates this audit trail. Backward/forward citation
chasing begins after an eligible full text is found; log seed,
direction, date, source and candidate IDs, and apply the same
eligibility and screening rules.

\hypertarget{6-publication-period}{%
\section{6. Publication period}\label{6-publication-period}}

Search for English-language works dated \textbf{1 January 2020 through
27 September 2026 inclusive} (date of draft). The final group protocol
must state its actual search cutoff and search date. If an interface
filters by year only, inspect 2026 records and exclude post-cutoff
publications manually. An older seminal mechanism or indispensable
baseline may enter only with a per-study written reason and a
demonstrated direct connection to the RQs; older contextual work is
otherwise background, not part of the included set. The assignment
targets 15--20 \textbf{unique included primary studies}, minimum 15, per
group; do not relax eligibility just to reach that count.

\hypertarget{7-inclusion-criteria}{%
\section{7. Inclusion criteria}\label{7-inclusion-criteria}}

All conditions must hold after full-text assessment:

\begin{enumerate}
\def\labelenumi{\arabic{enumi}.}
\tightlist
\item
  English, peer-reviewed journal or full conference paper reporting an
  original empirical or technical evaluation; ordinarily within §6.
\item
  A clearly identifiable local GPU sharing, partition,
  process-concurrency or API-interception/remoting mechanism involving
  at least two inference tenants, services, models or processes on a
  physical GPU. A transferable mixed workload qualifies only if
  inference-specific evidence and transfer rationale are explicit.
\item
  Sufficient full-text detail to identify workload, sharing
  configuration and at least one RQ-relevant engineering outcome or
  directly assessable isolation property.
\item
  A unique study, with multiple conference/journal/preprint reports
  linked and the most complete peer-reviewed report selected for
  extraction.
\end{enumerate}

Eligibility is independent of whether a mechanism performs well. Report
negative and null findings. Record any justified older exception and the
group\textquotesingle s adjudication.

\hypertarget{8-exclusion-criteria}{%
\section{8. Exclusion criteria}\label{8-exclusion-criteria}}

At full text record one primary reason, with secondary notes if useful:
\textbf{E1} no inference-serving evidence; \textbf{E2} no relevant
single-node multi-tenant sharing; \textbf{E3} mechanism outside assigned
scope; \textbf{E4} training-only; \textbf{E5} cluster orchestration
without local sharing evaluation; \textbf{E6} no peer-reviewed original
primary study (survey, vendor page, blog, preprint-only, poster,
editorial, patent); \textbf{E7} insufficient technical/outcome evidence;
\textbf{E8} duplicate/less complete report of the same study;
\textbf{E9} outside date/language rule without exception; \textbf{E10}
other, with written justification. Vendor guides {[}10{]}, {[}11{]}
describe mechanism context but cannot count as primary evidence or
demonstrate workload performance. Track full text \textbf{not retrieved}
as an access outcome, with attempts, rather than an exclusion among
assessed reports {[}2{]}.

\hypertarget{9-screening-procedure}{%
\section{9. Screening procedure}\label{9-screening-procedure}}

\begin{enumerate}
\def\labelenumi{\arabic{enumi}.}
\tightlist
\item
  \textbf{Identification:} Execute and export each frozen query. Assign
  a unique record ID to every hit, preserving database/run provenance.
  Log citation-chased candidates separately.
\item
  \textbf{Deduplication and versions:} Normalize DOI, title, year and
  authors; retain raw records and link record IDs → report IDs →
  underlying unique study IDs. Prefer the complete peer-reviewed version
  without treating its preprint as a second study.
\item
  \textbf{Calibration:} Both screeners independently apply the draft
  rule set to a small common pilot set; record disagreements and clarify
  ambiguous terms before freeze. The group\textquotesingle s third
  member adjudicates.
\item
  \textbf{Title/abstract:} Two independent screeners mark include,
  exclude or uncertain with reason and IDs. Retrieve uncertain records;
  adjudicate disagreement without inventing a consensus.
\item
  \textbf{Full text:} Record retrieval attempts and access status. Two
  screeners independently apply §§7--8 to retrieved reports; the third
  adjudicates discrepancies and records the final reason and
  study/report links. No inference from an unavailable full text.
\item
  \textbf{Flow:} Derive PRISMA 2020 record, report and study counts from
  those logs, including pre-screen removals, reports sought, reports not
  retrieved, assessed reports, exclusions with reasons and included
  unique studies {[}2{]}. Reconcile counts by group and by
  formal/supplementary discovery route.
\end{enumerate}

If the pilot suggests fewer than 15 eligible studies, consult the
instructor before changing the assigned scope or criteria. Do not
retrospectively turn scoping candidates into screened inclusions.

\hypertarget{10-quality-assessment-criteria}{%
\section{10. Quality-assessment
criteria}\label{10-quality-assessment-criteria}}

Two reviewers independently score each \textbf{included} study on
Q1--Q10: \textbf{1} adequate, \textbf{0.5} partial/unclear, \textbf{0}
absent/not reported. Reconcile with a recorded rationale; the third
member resolves persistent disagreement. This group-designed rubric is
an appraisal plan, not a validated universal instrument {[}4{]}.

\begin{longtable}[]{@{}p{0.23\textwidth}p{0.70\textwidth}@{}}
\toprule\noalign{}
ID & Question \\
\midrule\noalign{}
\endhead
\bottomrule\noalign{}
\endlastfoot
Q1 & Is the engineering question and tenant setting explicit? \\
Q2 & Are the mechanism and isolation/configuration described
sufficiently? \\
Q3 & Are GPU SKU/count, partition, OS, driver and runtime stated? \\
Q4 & Are model, workload, batching/arrivals and co-tenant mix stated? \\
Q5 & Is the baseline appropriate and configured fairly? \\
Q6 & Are utilization, latency percentiles and other metric
definitions/units clear? \\
Q7 & Are warm-up, duration, repetitions/variance and load generation
adequate? \\
Q8 & Do figures/tables/data support the reported findings? \\
Q9 & Are limitations, confounders and validity threats discussed? \\
Q10 & Is there enough information to reproduce or critically interpret
the comparison? \\
\end{longtable}

Total 0--10: \textbf{high 8--10}, \textbf{moderate 6--7.5}, \textbf{low
\textless6}. Score does not replace eligibility. Retain eligible
low-quality evidence with labels and rerun conclusions without it. Log
reviewer-specific scores, evidence anchors, reconciliation and missing
facts.

\hypertarget{11-data-extraction-fields}{%
\section{11. Data-extraction fields}\label{11-data-extraction-fields}}

Use one study-level identity row plus report/result rows where needed,
each with PDF page/section/table/figure anchors. Pilot the form on
multiple mechanism families before freeze.

\begin{longtable}[]{@{}p{0.23\textwidth}p{0.70\textwidth}@{}}
\toprule\noalign{}
Field group & Required fields \\
\midrule\noalign{}
\endhead
\bottomrule\noalign{}
\endlastfoot
Provenance & Group, record/report/study IDs, query/run and database,
DOI/publisher URL, authors, title, year, venue, peer-review/version
status, access/retrieval date, source anchor \\
Mechanism & Time-slice policy or context switch; MPS version/limits; MIG
profile; API remoting/proxy path; hybrid; scheduler/control layer,
admission/QoS policy \\
Environment & GPU vendor/SKU/count/memory/partition, host CPU/RAM,
OS/kernel, driver, CUDA/runtime, container or VM \\
Workload & Inference task/model/size/input, tenant identity and mix,
batch size, arrival pattern, concurrency, load,
warm-up/window/repetitions \\
Comparator & Baseline, direct vs indirect comparison, configuration
equivalence \\
Outcomes & Utilization/occupancy, throughput/goodput, tenant density,
p50/p95/p99 and units, SLO threshold/miss rate, interference, overhead,
energy if given, uncertainty/variance \\
Isolation & \textbf{Separately} measured or claimed performance
interference, resource/memory partitioning, security boundary and fault
containment; mark not reported where absent \\
Interpretation & Verbatim-sized finding summary in own words, exact
anchor, author caveat, reviewer inference separately, QA scores,
confounders and comparability notes \\
\end{longtable}

Do not fill absent values by assuming a GPU mechanism\textquotesingle s
default behavior. Capture both favorable and unfavorable findings.

\hypertarget{12-evidence-synthesis-strategy}{%
\section{12. Evidence-synthesis
strategy}\label{12-evidence-synthesis-strategy}}

Build a mechanism taxonomy (temporal, MPS, MIG, API
remoting/virtualization, hybrids), an RQ-to-study matrix and a
condition/comparability table. For RQ1 summarize evaluated control and
isolation dimensions; for RQ2 compare metric and experiment definitions;
for RQ3 stratify by GPU generation/profile, model and
batch/arrival/tenant conditions; for RQ4 trace reported trade-offs and
counterexamples to differing baselines and conditions; for RQ5 create a
conditional decision matrix with evidence strength and uncertainty.

Use structured narrative synthesis because different GPUs, workloads,
tenancy and latency percentiles make raw pooling misleading. Compare
numerical effects only inside genuinely compatible experimental strata;
otherwise report direction, conditions and limitations without an
invented aggregate. Report each recommendation as \textbf{direct
within-study comparison}, \textbf{indirect primary evidence},
\textbf{vendor context only}, or \textbf{insufficient evidence}. Apply
the SWiM transparency principles for any synthesis without meta-analysis
{[}9{]}. Repeat the conclusion after excluding low-QA/incomplete
studies. Report coverage gaps, publication/access bias, heterogeneity
and reviewer disagreement. No mechanism is declared universally optimal.

\hypertarget{protocol-changes-and-freeze-record}{%
\section{Protocol changes and freeze
record}\label{protocol-changes-and-freeze-record}}

\textbf{A001 (historical):} v1.0 was prematurely called frozen; v1.1
corrected the proposed databases, recall risk and full-text accounting
before any formal screening. Details remain in
\texttt{evidence/protocol\_amendments.md}; no count impact was claimed.

\textbf{A002 (this proposal, 27 September 2026 PKT):} Reorganize the
shared protocol into the instructor\textquotesingle s exact twelve
headings, expand methodological and mechanism references, specify
executable-query validation and group-specific ownership, and align
report source and Markdown. This is a \textbf{pre-search
reporting/clarification revision}, not a claim of completed search or
changed eligibility. Once the two groups are named, each must evaluate
these choices independently and record its own dated freeze; A002
affects zero formal records based on the current documented status.

\textbf{After a group starts formal searching:} Before using a major
change, append a dated entry to \emph{that group\textquotesingle s}
\texttt{evidence/protocol\_amendments.md} with old/new rule, trigger,
justification, approving reviewers, affected queries/dates/records,
whether earlier databases must be rerun, screening/PRISMA count effects,
and how old and new results will be reconciled. Version the protocol and
label deviations in the report. Log purely syntactic platform
translations in search logs; a recall-changing translation is a major
amendment. PRISMA-P calls for recording important protocol amendments
{[}1{]}.

\hypertarget{references-ieee-style-context-and-method-only-no-included-studies-claimed}{%
\section{References (IEEE style; context and method only, no included
studies
claimed)}\label{references-ieee-style-context-and-method-only-no-included-studies-claimed}}

{[}1{]} L. Shamseer \emph{et al}., ``Preferred reporting items for
systematic review and meta-analysis protocols (PRISMA-P) 2015:
elaboration and explanation,'' \emph{BMJ}, vol. 349, g7647, 2015. doi:
\href{https://doi.org/10.1136/bmj.g7647}{10.1136/bmj.g7647}.

{[}2{]} M. J. Page \emph{et al}., ``The PRISMA 2020 statement: an
updated guideline for reporting systematic reviews,'' \emph{BMJ}, vol.
372, n71, 2021. doi:
\href{https://doi.org/10.1136/bmj.n71}{10.1136/bmj.n71}.

{[}3{]} M. L. Rethlefsen \emph{et al}., ``PRISMA-S: an extension to the
PRISMA statement for reporting literature searches in systematic
reviews,'' \emph{Systematic Reviews}, vol. 10, art. 39, 2021. doi:
\href{https://doi.org/10.1186/s13643-020-01542-z}{10.1186/s13643-020-01542-z}.

{[}4{]} B. Kitchenham and S. Charters, ``Guidelines for performing
systematic literature reviews in software engineering,'' EBSE Technical
Report EBSE-2007-01, 2007. {[}Online{]}. Available:
\href{https://ebse.webspace.durham.ac.uk/ebse-bibliography/guidelines-for-performing-systematic-literature-reviews-in-software-engineering/}{EBSE
record}.

{[}5{]} Engr. Abul Kalam Library, NED University, ``E-Resources,''
accessed Sept. 27, 2026. {[}Online{]}. Available:
\href{https://eaklibrary.neduet.edu.pk/eakl/E-Resources.html}{NED
E-Resources}.

{[}6{]} IEEE, ``IEEE Xplore: Searching and Saving Searches,'' user
guide. {[}Online{]}. Available:
\href{https://ieeexplore.ieee.org/Xplorehelp/downloads/user-guides/IEEE_Xplore_Searching_and_Saving_Searches.pdf}{IEEE
guide}.

{[}7{]} ACM, ``Search tools,'' ACM Digital Library. {[}Online{]}.
Available:
\href{https://libraries.acm.org/training-resources/search-tools}{ACM
search tools}.

{[}8{]} Springer Nature, ``Springer Nature Link advanced search
option.'' {[}Online{]}. Available:
\href{https://support.springernature.com/en/support/solutions/articles/6000080445-springer-nature-link-advanced-search-option}{Search
help}.

{[}9{]} M. Campbell \emph{et al}., ``Synthesis without meta-analysis
(SWiM) in systematic reviews: reporting guideline,'' \emph{BMJ}, vol.
368, l6890, 2020. doi:
\href{https://doi.org/10.1136/bmj.l6890}{10.1136/bmj.l6890}.

{[}10{]} NVIDIA, ``Introduction,'' \emph{Multi-Instance GPU User Guide}.
{[}Online{]}. Available:
\href{https://docs.nvidia.com/datacenter/tesla/mig-user-guide/introduction.html}{MIG
introduction}.

{[}11{]} NVIDIA, ``Architecture,'' \emph{Multi-Process Service}.
{[}Online{]}. Available:
\href{https://docs.nvidia.com/deploy/mps/architecture.html}{MPS
architecture}.

{[}12{]} Z. Bai \emph{et al}., ``PipeSwitch: Fast pipelined context
switching for deep learning applications,'' in \emph{Proc. 14th USENIX
OSDI}, 2020. {[}Online{]}. Available:
\href{https://www.usenix.org/conference/osdi20/presentation/bai}{USENIX
venue page}. \textbf{Search seed only; not screened.}

{[}13{]} M. Lee \emph{et al}., ``ParvaGPU: Efficient spatial GPU sharing
for large-scale DNN inference in cloud environments,'' in \emph{Proc. SC
\textquotesingle24}, 2024. doi:
\href{https://doi.org/10.1109/SC41406.2024.00048}{10.1109/SC41406.2024.00048}.
\textbf{Search seed only; not screened.}

{[}14{]} M. Yu \emph{et al}., ``Torpor: GPU-enabled serverless computing
for low-latency, resource-efficient inference,'' in \emph{Proc. USENIX
ATC \textquotesingle25}, 2025. {[}Online{]}. Available:
\href{https://www.usenix.org/conference/atc25/presentation/yu}{USENIX
venue page}. \textbf{Search seed only; not screened.}

\end{document}
